Network operators must infer remaining path capacity from incomplete measurements, while avoiding measurement traffic that perturbs the path being measured. This project studies whether the packet-delay distribution observed at the receiver can provide a useful passive signal for this problem. Rather than reducing delay to a mean, percentile, or scalar counter, the monitor reconstructs an empirical delay distribution, applies adaptive stopping to decide when that reconstruction is stable, and then compares stopped distributions as offered load changes.

The statistical evaluation connects reconstruction to finite-domain distribution learning and uses total variation distance to measure how quickly discrete and continuous delay distributions can be learnt. The tested distributions converge well before the worst-case theoretical sample bound. For the continuous normal, lognormal, and Weibull cases, kernel density estimation reaches the 0.05 total variation distance target at 500 samples.

The online evaluation replaces access to ground truth with adaptive stopping. Consecutive cumulative reconstructions are compared after packet batches, and the experiments identify \(B=200\) and \(\theta=\delta/6\) as a reliable continuous operating point across the three tested delay families. This setting is treated as an empirical operating choice, not as a theoretical guarantee.

The capacity use case is evaluated in Network Simulator 3 under constant-rate cross traffic, Pareto on-off bursty traffic, and a heterogeneous slow-link stress test. The initial hypothesis that first distributional change marks capacity is rejected because it fires too early. A stronger criterion based on the rate of change of total variation distance detects a capacity-relevant boundary at the 0.925 load-grid point in both single-link and two-to-ten-hop constant-rate experiments. Under bursty traffic, the detected average-load boundary moves to 0.475 because the on-period traffic rate is twice the swept average. In the heterogeneous five-hop test, the same 0.925 grid point is recovered when the slow link is first, middle, or final. A final constructed multi-path modal experiment shows that, in the tested separable mixture, the capacity response should be associated with the changing mode rather than the whole mixture. The work does not provide an exact physical capacity measurement. It shows that passive distributional monitoring can expose when a path leaves its previous low-load operating regime.
